\chapter{The Rationality of Everyday Causal Cognition}
\label{sec:p19-rationality-7-2}

Waldmann's question is not simply whether people can identify a causal
effect from a probability distribution. It is which representations and
learning procedures people use, which norm is appropriate for a given
task, and how their judgments fare under that norm. His selective review
connects competing philosophical accounts with competing psychological
accounts, rather than declaring all causal cognition either rational or
irrational. Its main distinction is between dependency representations
of variables and dispositional representations of interacting objects.
Neither a graph-theoretic identification result nor an isolated behavioral
success supplies the whole treatment
\cite[chapter~7.2, pp.~439--448]{Waldmann2021EverydayCausal}.

The empirical scope is correspondingly bounded. The chapter supplies
the conceptual argument and its attributed review-level findings. The
two experiments of Meder, Hagmayer and Waldmann supply the detailed
trial-by-trial learning record below
\cite{MederHagmayerWaldmann2008Interventional}.
No pooled effect, participant-level reanalysis, universal accuracy rate
or exhaustive replication claim is inferred from this selection.

\section{Which question is being evaluated?}
\label{sec:p19-rat72-tasks}

\begin{hbdefinition}{def:p19-rat72-task-distinction}{Causal questions with different targets}
A \emph{predictive} question asks about effects given information about
causes; a \emph{diagnostic} question asks about causes given information
about effects. A \emph{strength} question concerns a specified causal
relation's parameter, whereas a \emph{support} question concerns whether
such a relation exists. A \emph{structure-learning} question concerns
which causal relations connect the variables. An \emph{interventional}
question asks what happens under a specified manipulation, not merely
upon observing its target value. A \emph{singular-causation} question
asks whether one particular occurrence produced another occurrence.
Finally, choosing a causal verb or assigning agent/patient roles can
concern how an interaction is represented linguistically or perceptually.
These descriptions identify tasks, not interchangeable numerical scores.
\end{hbdefinition}

Knowing that smoking raises the frequency of lung disease does not by
itself answer the chapter's question whether Peter's smoking caused
Peter's disease. Nor does successfully predicting weather from a
barometer establish that changing the barometer changes the weather.
A participant may learn the frequency of an effect, infer a causal
structure, or interpret the examiner's word ``strength'' as evidence
that a link exists. Those possibilities differ even when the response
is one number. The task, stimuli, information and response wording
must therefore precede the choice of benchmark, as in
\cref{def:p19-empirical-question,def:p19-benchmark}.

Waldmann emphasizes a constitutive connection between normative and
descriptive inquiry. Scientists and lay reasoners both ordinarily mean
their causal claims to be correct. Psychological models consequently
often borrow from philosophical or methodological models: association
from an empiricist tradition, causal graphs from probabilistic and
causal modeling, and object-based powers from dispositional accounts.
This shared aim does not make the resulting methods or representations
identical. A procedure developed for economic variables need not be an
adequate account of intuitive collisions, and a useful local rule need
not be part of a globally coherent causal theory. His historical account
also attributes the earlier relative neglect of causal reasoning in
mainstream cognitive psychology to domain-general learning and
similarity-based approaches. That is his explanatory historical claim,
not a count of all textbooks or a new survey here
\cite[Secs.~1--2, pp.~439--440]{Waldmann2021EverydayCausal}.

\begin{hbconvention}{def:p19-rat72-evaluation}{Two standards and three kinds of claims}
\emph{Coherence} evaluates conformity to a specified normative model.
\emph{Correspondence} evaluates whether representations or predictions
match the relevant relations in the world. A formal implication is
conditional on its model's assumptions; a descriptive finding concerns
recorded responses in a task and sample; a normative prescription needs
an argument for applying its standard to that task. A close descriptive
fit alone establishes neither correspondence nor a prescription.
\end{hbconvention}

There is no uncontested single norm of causation in the review. A response
that violates a simple covariation norm may accord with inference over
uncertain causal models, while agreement with that inference can still
depend on erroneous prior beliefs. Some theories respond to failures by
improving descriptive fit without retaining the original normative claim;
others revise the normative model, for example by admitting uncertainty
previously ignored. These are different repairs. The latter must justify
the additional assumptions rather than treating every discrepancy as
evidence for an unspecified hidden variable. The source also reports
distortions in encoding probabilities and especially zero covariation;
prior knowledge may explain some but not all. Its conjecture that
ordinary causal knowledge is often adequate because it is adaptive is
not a measured population accuracy rate
\cite[Sec.~2, p.~440]{Waldmann2021EverydayCausal}.

\section{Dependency representations: direction, strength and support}
\label{sec:p19-rat72-dependency}

The dependency framework represents causal relations using variables
and their possible values. Its members include associative, covariation,
power-PC, causal-model and Bayesian-inference theories; they are not
one algorithm. Their causal relata can be events, propositions, facts,
properties or states of affairs under different philosophical accounts.
The source's forest-fire example uses the presence/absence of a lit
match and lightning as two potential causes of a fire. Its Figure~7.2.1
distinguishes a common effect, with two cause-to-effect arrows, from a
common cause, with two effect variables downstream. Mechanistic
knowledge can enter such representations through mediating variables.
Representing a relation this way is not the unrestricted principle that
any marginal statistical dependence proves a causal arrow
\cite[Sec.~3.1, p.~441]{Waldmann2021EverydayCausal}.

\subsection{Presentation order is not causal direction}

In an associative learning account, encountered cues acquire weights
predicting subsequently supplied outcomes. A symptom may be encountered
first and a disease diagnosis second, even though the disease causes
the symptom. A causal-model account instead represents the direction
of production independently of the order in which information arrives.
Thus the two accounts can agree about a cue's predictive success while
disagreeing about which other cues should compete with it and which
manipulations should propagate through the system.

Waldmann's review attributes direction-sensitive cue-competition
evidence to Waldmann and Holyoak~(1992), Waldmann~(2000) and related
work. The learning contrast concerns predictive versus diagnostic
arrangements, not a survey of physicians' accuracy. The review also
attributes fallback to simpler associations in complex learning tasks
to Lopez, Cobos and Cano~(2005). These are qualitative review records:
the chapter does not give their recruitment frames, individual sample
counts, assignment procedures, numerical measures or sampling intervals.
Consequently they support the chapter's report of a debate and its
complexity qualification, not a numerical prevalence estimate or the
claim that every learner represents a graph. A flexible associative
account and a causal representation with limited performance remain
alternatives unless their task-specific predictions are separated
\cite[Sec.~3.1, p.~441]{Waldmann2021EverydayCausal}.

\subsection{Covariation, power and evidence for a link}

For binary variables \(C,E\), let \(q_1=P(E=1\mid C=1)\) and
\(q_0=P(E=1\mid C=0)\), with both conditioning events of positive
probability. The difference \(\Delta P=q_1-q_0\) measures the observed
change in effect frequency. It does not name every possible causal
strength construct. The source contrasts it with Cheng's causal power:
the theoretical probability that a cause, acting alone, generates or
prevents an effect. Because the target is not actually observed in
isolation, the inference requires background assumptions, including
appropriate independence from other causes. Ceiling effects are decisive:
if another cause already guarantees the effect, an additional generative
cause cannot further raise its observed frequency. A zero frequency
contrast then need not mean a powerless mechanism
\cite[pp.~441--442]{Waldmann2021EverydayCausal}.

\begin{hbtheorem}{thm:p19-rat72-power-benchmark}{proved=annex:p19.rat72-power-benchmark}{A bounded generative-power comparison}
Let \(C,K,T\) be mutually independent binary random variables, with
\(0<P(C=1)<1\), \(P(K=1)=r\), \(P(T=1)=s\), and
\(E=K\lor(C\land T)\). Then
\[
 q_0=r,\qquad q_1=r+(1-r)s,\qquad \Delta P=(1-r)s.
\]
If \(r<1\), \(s=\Delta P/(1-q_0)\); if \(r=1\), the two observed
conditional probabilities equal one for every \(s\), so they do not
determine the target's success probability. These are implications of
this stipulated independent-generators model, not empirical laws.
\end{hbtheorem}

For example, the same contrast \(0.25\) arises with \(r=0,s=0.25\)
and with \(r=0.5,s=0.5\). They differ in what the target can produce
without background generation. Likewise, adding two independent
generative strengths \(0.75+0.75\) would exceed one; the probability
that at least one generator succeeds is instead
\(1-(1-0.75)^2=0.9375\). These are editorial illustrations of the
source's ceiling objection, not data or an assertion that all physical
causes combine by noisy-OR. Preventive causes and dependent background
mechanisms need their own models.

The next alternative changes the target rather than just rescaling
\(\Delta P\). Griffiths and Tenenbaum's causal-support account, as
reconstructed by Waldmann, compares a model containing a \(C\)-to-\(E\)
link with a model in which background causes alone influence \(E\).
It handles parameter uncertainty by distributions rather than one fitted
power value. In the elementary Bayesian comparison, posterior odds for
a link equal prior odds multiplied by the likelihood ratio after each
model's unknown parameters have been integrated. Such a posterior can
remain positive after an observed zero contingency; zero empirical
contrast is not certainty that the no-link model is true. Whether a
participant is being asked for power, observed contrast, or support
for the existence of a link therefore matters. The claim that support
is what learners intend by a strength response is a psychological
hypothesis, not a consequence of Bayes' rule
\cite[p.~442]{Waldmann2021EverydayCausal}.

\section{Seeing and doing after observational learning}
\label{sec:p19-rat72-intervention}

In the chapter's barometer example, atmospheric pressure is a common
cause of both the reading and weather. Observing the reading gives
information about pressure and hence weather. A manipulation that
sets only the reading does not change pressure or weather. Causal-model
reasoning distinguishes these operations by removing the reading's
ordinary incoming causal relation under the intervention. Pure
cue--outcome associations, without further representational machinery
or prior instrumental experience, do not supply this transformation.
The conclusion nevertheless requires the common-cause hypothesis and
invariance of the untouched mechanisms. Observational frequencies
alone do not license graph surgery on an arbitrary, unidentified system
\cite[p.~442]{Waldmann2021EverydayCausal}.

\subsection{The four-variable learning task and its benchmark}

Meder, Hagmayer and Waldmann tested whether the distinction survives
trial-by-trial learning rather than being restricted to tabulated data.
Their instructed graph is the diamond
\(A\to B\to D\) and \(A\to C\to D\), using four fictitious chemical
substances in wine casks. The diagram initially supplies the structure,
not the numerical parameters, and is absent during learning and testing.
Forty passively observed cases display presence or absence; participants
do not practise the queried interventions. The study therefore tests
inference given instructed structure and learned parameters, not
unassisted discovery of the graph
\cite[pp.~75--77]{MederHagmayerWaldmann2008Interventional}.

\begin{hbconvention}{def:p19-rat72-diamond-notation}{The experimental finite comparison}
The capitals \(A,B,C,D\) now denote binary random variables in one
stipulated graph, and lower-case indices \(a,b,c,d\) denote their
values. Set \(r=P(A=1)\), with \(0<r<1\),
\(u_a=P(B=1\mid A=a)\), and \(v_a=P(C=1\mid A=a)\).
Write \(t_{bc}\) for the specified mechanism probability of \(D=1\)
at parent values \(B=b,C=c\). It agrees with
\(P(D=1\mid B=b,C=c)\) on positive-probability parent cells;
the mechanism remains specified even on unobserved parent cells.
Assume a finite structural causal model on this graph with mutually
independent disturbances and these conditional kernels. Its joint law
then factors as
\[
 P(a,b,c,d)=P(a)P(b\mid a)P(c\mid a)P(d\mid b,c).
\]
The last factor denotes the specified normalized mechanism kernel,
including its values on zero-probability parent configurations.
The notation \(\operatorname{do}(C=c)\) means fixed replacement of
the equation for \(C\), retaining the other mechanisms and exogenous
law, as in \cref{def:p16-causal-intervention}. The structural and
invariance assumptions are additional causal content, not consequences
of a factored observational table alone. It is not conditioning.
\end{hbconvention}

\begin{hbtheorem}{thm:p19-rat72-diamond-benchmark}{proved=annex:p19.rat72-diamond-benchmark}{Diagnostic and predictive benchmarks in the diamond}
For \(c=0,1\), put \(\ell_a(c)=P(C=c\mid A=a)\), and suppose
\(P(C=c)=r\ell_1(c)+(1-r)\ell_0(c)>0\). Define
\[
 k_c=\frac{r\ell_1(c)}{r\ell_1(c)+(1-r)\ell_0(c)},\quad
 h_c=k_cu_1+(1-k_c)u_0,\quad h=ru_1+(1-r)u_0.
\]
Then
\[
\begin{aligned}
 P(A=1\mid C=c)&=k_c,&
 P(A=1\mid\operatorname{do}(C=c))&=r,\\
 P(D=1\mid C=c)&=t_{1c}h_c+t_{0c}(1-h_c),&
 P(D=1\mid\operatorname{do}(C=c))&=t_{1c}h+t_{0c}(1-h).
\end{aligned}
\]
The first predictive expression uses the observationally updated root;
the second preserves the root's base rate. This fully indexed
reconstruction supplies the intended computations in the source's
Appendix~A1--A8.
\end{hbtheorem}

The printed Figure~1 parameters are \(r=0.5\), \(u_1=0.85\),
\(u_0=0.10\), \(v_1=0.90\), \(v_0=0.05\),
\(t_{00}=0\), \(t_{10}=t_{01}=0.75\), \(t_{11}=0.93\).
Using that printed two-causes value as a finite calibration gives
\(k_1=18/19\), \(k_0=2/21\), \(h_1=77/95\),
\(h_0=6/35\), and \(h=0.475\). Thus observing absent \(C\)
suggests absent \(A\) and often absent \(B\); actively preventing
\(C\) leaves the \(A\to B\to D\) route available. The predicted
probability of \(D\) is about \(0.129\) under the former operation
and \(0.356\) under the latter. The source reports its benchmark
as whole-number frequencies out of forty. Those rounded frequencies
are not measured participant responses
\cite[Fig.~1, Tables~1--2 and Appendix, pp.~76--80]
{MederHagmayerWaldmann2008Interventional}.

\subsection{Design, measures and the separating results}

Experiment~1 recruited 24 University of Goettingen students for course
credit. Observation/intervention and present/absent \(C\) were
within-person factors. Each trial showed \(A\), then \(B,C\), then
\(D\), with one-second interstimulus intervals and a further two-second
display of the completed pattern. The eight test questions asked for
conditional frequencies of \(A\) or \(D\) in forty new cases, under
either observation or addition/prevention of \(C\). Question types
were blocked, with block order counterbalanced. Experiment~2 used
24 other Goettingen participants and the same data and graph, but
displayed \(D\), then \(B,C\), then \(A\). The causal model did not
reverse when the learning presentation reversed.

\begin{center}
\small
\begin{tabular}{lrrr}
\hline
Query & Benchmark & Experiment~1 & Experiment~2\\
 & count/40 & mean (SD) & mean (SD)\\
\hline
\(A\mid C=1\) &38&30.5 (7.6)&33.5 (8.6)\\
\(A\mid C=0\) &4&17.1 (10.4)&15.4 (11.2)\\
\(A\mid\operatorname{do}(C=1)\)&20&25.5 (10.6)&25.5 (11.2)\\
\(A\mid\operatorname{do}(C=0)\)&20&27.3 (8.6)&22.4 (10.3)\\
\(D\mid C=1\)&36&29.7 (10.0)&30.5 (10.1)\\
\(D\mid C=0\)&5&14.8 (11.6)&18.3 (13.3)\\
\(D\mid\operatorname{do}(C=1)\)&33&27.5 (11.6)&29.7 (10.9)\\
\(D\mid\operatorname{do}(C=0)\)&14&21.6 (12.6)&20.3 (11.7)\\
\hline
\end{tabular}
\end{center}
The SDs describe between-participant dispersion of the recorded frequency
estimates, not confidence intervals for population means. The paper
does not supply those intervals for these entries
\cite[Table~2, p.~77]{MederHagmayerWaldmann2008Interventional}.

In Experiment~1 the diagnostic observational contrast was significant,
\(F(1,23)=36.51,p<0.001\), while the two interventional diagnostic
estimates did not differ significantly, \(F<1\). The key predictive
comparison, absent observed \(C\) versus prevented \(C\), was also
significant, \(F(1,23)=9.57,p<0.01\). Crucially, present observed
versus generated \(C\) did not yield a significant predictive
difference, \(F(1,23)=1.0,p=0.33\), which constrains a simple
``always answer doing differently'' response explanation. It is not
an equivalence test.

Experiment~2 retained diagnostic differentiation: the observed present/
absent contrast was \(F(1,23)=63.88,p<0.001\); the interventional
contrast was \(F(1,23)=1.52,p=0.23\). But the observation-versus-
intervention predictive interaction and the absent/prevented predictive
comparison both had \(F<1\). Thus a direction-sensitive diagnostic
success coexists with failure to demonstrate the more demanding
alternative-pathway contrast under reversed learning order. Comparing
one significant and one nonsignificant result does not itself prove a
between-experiment interaction. The authors propose imperfect parameter
learning as an explanation: in Experiment~2 participants' separately
elicited mean estimate of \(P(A=1)\) was \(0.65\), rather than the
displayed \(0.5\). That observation is consistent with, not proof of,
their proposed mechanism
\cite[pp.~77--79]{MederHagmayerWaldmann2008Interventional}.

These results support a bounded capacity to distinguish observation
and intervention without practising the relevant actions. They do not
show exact numerical calibration, population representativeness, one
unique mental algorithm, or reliable structure discovery. Instructed
graph knowledge, frequency memory, interpretation of the hypothetical
questions and approximate mental simulation can each contribute.
The paper itself distinguishes competence from performance and leaves
the process of parameter learning open. Its unpublished follow-ups
mentioned in the discussion are not independently verified replications
here. The source chapter's further reports about human studies and
preliminary rat studies retain their different tasks and populations;
the student result cannot be transferred to animals or daily decisions
\cite[p.~442]{Waldmann2021EverydayCausal}
\cite[pp.~78--79]{MederHagmayerWaldmann2008Interventional}.

\section{Uncertainty, acquisition and apparent violations}
\label{sec:p19-rat72-acquisition}

\subsection{Diagnostic queries and singular cases}

The chapter uses Meder, Mayrhofer and Waldmann~(2014) to illustrate a
change in the normative account itself. Participants' diagnostic
judgments were lower when the predictive causal strength was low,
instead of simply reporting an observed conditional frequency of cause
given effect. A Bayesian alternative averages over uncertainty about
the causal model. For a specified family of models, its target is a
posterior-weighted average of the models' diagnostic predictions, not
the empirical frequency in isolation. There is no contradiction in
obtaining different predictions when the models or their posterior
weights differ despite equal observed diagnostic frequencies.

This is Waldmann's attributed account of the study and model rationale.
His chapter does not supply the study's sample counts, recruitment,
assignment, numerical responses, error bars or model-fit uncertainty.
The evidence record therefore concerns a described diagnostic task and
judgment pattern, not a numerical effect estimate from that summary.
The separate original-study treatment in
\cref{sec:p19-rationality-7-3} supplies the diagnostic tasks, recruited
and analyzed samples, measures, results, uncertainty and model comparisons.
Those details are not ascribed to Waldmann's 2021 summary, and neither
a summary nor a close model fit establishes that the proposed uncertainty
representation is uniquely used.
Simpler learned frequency distortions or alternative interpretations
need task-specific tests. The separate question whether Peter's
smoking produced his particular illness is introduced as a singular
query, with competing causal versus coincidental explanations. The
chapter refers to probabilistic accounts including timing and preemption;
that motivating example is not an observed patient study
\cite[p.~443]{Waldmann2021EverydayCausal}.

\subsection{Learning a graph versus learning within a graph}

Machine-learning algorithms may infer candidate causal networks from
covariation under strong constraints, but that success does not show
that people estimate the same conditional independences or process the
same volume of data. Waldmann reports studies using small candidate
sets and typically three variables, including Steyvers and colleagues'
observational/interventional learning work. He characterizes performance
as generally poor unless conditions are favorable, with some benefit
from interventions. The chapter does not report a common denominator,
sample composition, accuracy scale or interval for these heterogeneous
studies; ``poor'' is its review characterization, not a population rate.
Restricting candidates, giving graph instructions, or permitting actions
changes the task, so success in the wine-cask experiment is not contrary
evidence about unassisted graph induction
\cite[p.~443]{Waldmann2021EverydayCausal}.

Knowledge-based induction offers another route. Temporal order,
interventions, explicit instructions and prior domain knowledge guide
initial structure hypotheses; covariation then informs parameters or
competing hypotheses. Hierarchical Bayesian accounts represent inference
at data, model and theory levels, allowing higher-level physical beliefs
to constrain the candidate models. The conceptual advantage is a
restricted search space, not a theorem that the priors are true.
The source's additional examples include different integration rules
for physical quantities, expected cause--effect latencies, and
probabilistic intuitive-Newtonian models. Their domain information is
part of the inference rather than noise to be removed automatically.
It may assist learning or systematically mislead it. Those examples
retain the chapter's cited research boundary; no newly quantified
learning advantage or universality across domains is asserted
\cite[pp.~443--444]{Waldmann2021EverydayCausal}.

\subsection{Collateral effects and the Markov comparison}

In a common-cause graph with independent mechanisms, knowing the cause
screens one effect off from the other effects. Waldmann reports Rehder
and Burnett's finding that people estimating an effect given its cause
were nevertheless influenced by the presence or absence of collateral
effects. That is a violation of the \emph{specified simple model's}
prediction. The reported task is conditional feature inference; the
chapter supplies neither participant counts and recruitment nor
numerical changes and uncertainty. It must not be converted into an
estimate of the frequency of Markov violations among all people.

One response is to reject the adequacy of that causal representation,
an objection Waldmann associates with Cartwright. Another is to enrich
the model with a shared unobserved mechanism, disabler, or object
disposition: collateral effects may then reveal whether the cause is
operating effectively. Marginal dependence can remain even when the
enlarged graph satisfies its own Markov condition. Associative-process
hybrids supply further alternatives. The enlarged model needs
independently motivated mechanism beliefs or new predictions; arbitrary
post hoc additions would not discriminate a theory. A formal proof of
screening off and an empirical claim about the graph a person uses are
therefore different claims
\cite[p.~444]{Waldmann2021EverydayCausal}.

\section{Dispositions, forces and causal language}
\label{sec:p19-rat72-disposition}

The disposition framework asks why an interaction occurs by considering
the powers of its participating objects. Its primary entities are
objects, not just variables recording event presence. Solubility can
be a relatively stable disposition; exertion of force can be temporary.
A suitable context permits a power to be expressed. This is a deeper
explanatory proposal, but it is not automatically the same mechanism
as a chain of mediating variables. The source acknowledges both the
historical Aristotelian ancestry and disagreements about whether
modern dispositional accounts improve scientific methodology. The
ancestry alone validates neither modern physics nor a psychological
model
\cite[Sec.~3.2, p.~444]{Waldmann2021EverydayCausal}.

\subsection{Cause, prevent, enable and despite}

Force-dynamic accounts distinguish an agent or affector from a patient
and represent their tendencies and opposing forces relative to a
possible end state. These can be physical forces or metaphorical
extensions to social influence. In the chapter's wind/man/bench
example, a man has no tendency to move toward the bench, the wind
acts against his tendency, and the resultant points toward the bench;
this configuration is assigned the term \emph{cause}. Related
configurations underlie \emph{prevent}, \emph{enable} and
\emph{despite}. This concerns causal-language meaning and perceived
tendencies, not a fitted event-frequency parameter.

The source attributes visual-scene and causal-term-choice tests to
Wolff and colleagues. The task displays stationary or moving people
or objects and asks participants to select causal expressions. It does
not provide their population frame, counts, condition-assignment details,
accuracy values or uncertainty. Thus it reports a testing strategy
and selected qualitative fit, not a complete effect estimate or a
universal account of causal vocabulary. A verbal selection can also
depend on pragmatic conventions and framing; fitting it does not
alone show a Newtonian force calculation in the participant.

The scope extends beyond one positive collision. Knocking away a pole
that keeps a tent upright removes a preventer and thereby permits the
tent to fall. Omission cases examine the forces left active when an
action is not taken. These examples explain what the extended framework
is intended to capture; they are not observed experiments or a general
proof that every absence is a cause. A chosen graph can also encode
prevention, but translating away the object dispositions or lexical
contrasts would lose the question this source asks
\cite[pp.~444--445]{Waldmann2021EverydayCausal}.

\subsection{Launching and causal asymmetry}

In the Michotte-style scene described in the chapter, moving ball
\(X\) reaches resting ball \(Y\); \(X\) stops and \(Y\) starts moving.
Waldmann reports White's findings that participants favor describing
\(X\) as launching \(Y\), rather than \(Y\) as stopping \(X\), and
estimate a greater force for \(X\). These are two measures, verbal
agent assignment and force judgments, not one causal-strength score.
The animations are two-dimensional visual displays. The chapter does
not report participant counts, recruitment, assignment details,
numerical force differences or intervals. Consequently it warrants
an attributed directional finding with unknown sampling uncertainty,
not an accuracy percentage or a universal perceptual bias.

For physical contact forces the Newtonian comparison is equality in
magnitude and opposite direction. Unequal estimates of those contact
forces conflict with that benchmark. But the physical symmetry alone
does not imply equal linguistic appropriateness of every description:
initial motion, attention, the question and conversational goals can
make one causal description more useful. Thus White's sensorimotor
explanation---stored experiences of acting on a resisting patient---is
a proposed process account, not established by Newton's law or by
the preferred wording alone. The chapter describes follow-up work
on prior context and perceived forces, preserving a live comparison
between dispositions, physical knowledge and task interpretation
\cite[p.~445]{Waldmann2021EverydayCausal}.

\section{Relations, use cases and the evidence boundary}
\label{sec:p19-rat72-use-and-boundary}

The dependency treatment connects to
\cref{sec:p16-finite-causal-models}: that chapter supplies equations,
independent-disturbance hypotheses, intervention invariance, and proofs
of observational/interventional differences. Those assumptions make
the wine-cask benchmark precise; they do not explain whether learners
encode the parameters accurately. The probability-net comparison in
\cref{sec:p16-rat42-comparisons-and-use} likewise distinguishes a
joint distribution from causal interpretation. A Bayesian calculation
under a selected model, a claim that the model is identified, and a
claim that people mentally represent that model require different
justifications. This chapter supplies the latter questions and their
task boundaries, rather than duplicating formal identification coverage.

\Cref{sec:p19-evidence-methods} gives the empirical record needed to
evaluate such claims. Here the two directly inspected experiments
specify sample, displayed graph/data, temporal order, response format,
means/SDs, planned comparisons and limits. The remaining chapter-level
review findings explicitly retain unavailable population or uncertainty
information. A task-specific qualitative success cannot fill those
fields, and a reported nonsignificant result is not an established null.

\begin{hbexample}{def:p19-rat72-lookup}{Successful selection and failure cases}
To ask whether a person can reason about an unpractised intervention
from observations, give the causal structure, specify the manipulation
and test both diagnostic and predictive questions. The wine-cask
results separate a diagnostic success from a more demanding predictive
failure under reversed presentation. Scoring only the diagnostic items
and declaring general mastery would fail the intended comparison.

To compare a frequency contrast with a generative parameter, first
specify the background-generation assumptions. The two equal-contrast
models above require different target powers. When the background
guarantees the effect, the observed zero contrast cannot identify the
target power. Applying that model to dependent or preventive mechanisms
without revision would fail.

To test a Markov prediction, state the assumed variables and mechanisms
before interpreting collateral-effect sensitivity. A shared hidden
disabler changes the prediction and requires independent justification.
Explaining every mismatch by an unspecified disabler would fail as a
discriminating cognitive account.

To analyze \emph{prevent} or a launching description, record the
objects, tendencies, resultant and response wording. Replacing that
semantic task with a causal-effect identification exercise would answer
a different question. Conversely, object-role judgments do not supply
an intervention distribution for an arbitrary many-variable system.
\end{hbexample}

Waldmann's conclusion is deliberately unresolved. Dependency models
are often applied to predictions in complex causal networks;
dispositional models to language and small visual interactions. Their
different task selections help explain why much competition occurs
within each framework rather than between them. Attempts at a single
unified theory have not displaced this division in the reviewed work.
Causal pluralism instead allows different models for different tasks;
hybrid accounts ask how their representations interact. The chapter's
specific hybrid example has dispositions conveyed by verbal instructions
shaping the causal graph used for subsequent inference. It is not
merely a list of theories, nor proof that this combination uniquely
accounts for causal cognition. The unitary, pluralistic and hybrid
alternatives remain competing research programs
\cite[Sec.~4, pp.~445--446]{Waldmann2021EverydayCausal}.

No theorem that humans are universally rational is promised or proved.
The source's full argument concerns the relation between normative
choices, representations, tasks and observed departures. The two new
finite results have their reciprocal proofs in the annex; they establish
the displayed benchmark under stated assumptions, not empirical adequacy
of an inference algorithm. The chapter's references preserve attribution
and further-reading boundaries, not a claim that all cited books and
experimental programs have been independently reproduced.

\index{causal cognition!everyday rationality}
\index{causal learning!observation and intervention}
\index{causal power}\index{causal support}
\index{causal cognition!dispositions and force dynamics}
